\chapter{Záver}
\label{Zaver}

Práca si kládla za cieľ navrhnúť a~implementovať riadiaci systém
pre interaktívnu múzejnú expozíciu, ktorý by zvládol synchronizáciu
fyzických efektov, multimédií a~vetvenie scenára podľa vstupu
návštevníka -- a~to bez priemyselných PLC a~komerčných show-control
licencií. Systém bol nasadený na reálnej inštalácii CHRONOS
a~overený v~podmienkach každodennej múzejnej prevádzky.

\section*{Dosiahnuté výsledky}

Na softvérovej úrovni tvorí jadro riešenia konečný stavový automat
s~externou JSON konfiguráciou. Scenár expozície -- časovanie, poradie
efektov, podmienky vetvenia -- je možné meniť editáciou textového
súboru bez akéhokoľvek zásahu do zdrojového kódu. Toto oddelenie
obsahovej vrstvy od implementácie enginu je v~praxi kľúčové:
múzeum môže upravovať scenáre samostatne, bez nutnosti volať
vývojára. Komunikácia medzi centrálnym Raspberry~Pi a~uzlami ESP32
cez protokol MQTT dosiahla počas testovania 100\,\% doručiteľnosť
príkazov pri priemernej latencii spätnej väzby 111{,}7\,ms.
Systém bežal niekoľko dní bez reštartu a~po výpadku napájania
sa obnovil do plne funkčného stavu v~priemere za 38 sekúnd.

Na hardvérovej úrovni sa rozhodnutie stavať na certifikovaných
priemyselných moduloch namiesto vlastných dosiek osvedčilo.
Počas testovania nevznikla žiadna porucha spôsobená hardvérovým
zlyhaním a~výmena komponentu v~prípade poruchy nevyžaduje znalosti
elektroniky -- len náhradu kus za kus.

\section*{Obmedzenia}

Systém má dve identifikované obmedzenia. Bezdrôtová Wi-Fi komunikácia
vyžaduje vyhradenú sieť oddelenú od verejného pripojenia múzea;
pri zdieľaní siete s~návštevníkmi hrozí degradácia odozvy. Druhým
obmedzením je prvotná konfigurácia nových uzlov: aktualizácia
firmvéru prebieha bezdrôtovo, avšak prvé nastavenie sieťových
parametrov si vyžaduje fyzický prístup k~zariadeniu. Pre jednorázové
nasadenie ide o~zanedbateľnú záťaž, pri škálovaní na desiatky
uzlov by však stálo za to tento krok automatizovať.

\section*{Výhľad do budúcna}

Najväčší priestor na rozvoj vidím v~rozšírení stavového automatu
o~perzistentné dáta v~rámci scény. Aktuálna implementácia
neuchováva stav medzi prechodmi; pridanie jednoduchého kľúč-hodnota
slovníka by otvorilo celú triedu nových scenárov -- napríklad
kvízové expozície, kde sa v~priebehu scény zbierajú body
a~záverečný stav (správna alebo nesprávna odpoveď, výsledné
skóre) určuje, do ktorého stavu automat prejde. Súvisiacim
rozšírením sú náhodné prechody: namiesto deterministickej
sekvencie by systém pri každom spustení mohol vybrať inú vetvu
scenára, čím sa predíde stereotypu pri opakovaných návštevách.

Druhým smerom je vizuálny editor scenárov, ktorý by nahradil
manuálne editovanie JSON súborov grafickým rozhraním. Scenár
by sa skladal z~blokov a~šípok priamo v~prehliadači -- bez
znalosti formátu konfiguračného súboru. Rozšírenie systému
na viacero miestností je z~architektonického hľadiska
pripravené: vyžaduje len pridanie nových uzlov a~konfiguráciu
príslušných topikov, nie zmenu kódu enginu.